\section{The Bombieri--Vinogradov small-divisor interface}
\label{sec:small-divisor-interface}

This section separates three statements which are easy to conflate.  First,
expanding the first von Mangoldt weight gives an exact divisor ledger; a
\emph{source divisor} means a divisor introduced by that expansion.  Second,
the Bombieri--Vinogradov theorem evaluates the part with
$d\le X^{1/2}(\log X)^{-B}$.  Third, the complementary large-divisor tail is
still unestimated.  We also record a different one-sided application
of Bombieri--Vinogradov in which one factor and the target are prime-weighted.
None of these statements supplies a lower bound for a fixed prime pair.

Throughout the first part of the section, $h\ne0$ is a fixed even integer,
$W\in C_c^1((0,\infty))$ is supported in $[1,2]$, and
\begin{equation}\label{eq:weight-mass}
 \cI_W=\int_0^\infty W(u)\,du.
\end{equation}
The restriction to even $h$ is the admissible two-form case.  The same
identities hold for odd $h$, but then the local singular series vanishes.

\subsection{An exact source-divisor ledger}
\label{subsec:source-divisor-ledger}

Consider the smoothed fixed-shift correlation
\begin{equation}\label{eq:pair-correlation}
 \cC_h(X;W)
 =\sum_{n\geq1}\Lambda(n)\Lambda(n+h)W(n/X),
\end{equation}
where $X$ is sufficiently large that every target in the support is positive.
The elementary identity
\begin{equation}\label{eq:mobius-log-source-identity}
 \Lambda(n)=-\sum_{d\mid n}\mu(d)\log d
\end{equation}
gives the following finite decomposition.

\begin{proposition}[Exact source M\"obius--log ledger]
\label{prop:exact-source-ledger}
For every $X$ and every $1\leq D\leq2X$,
\begin{align}
 \cC_h(X;W)
 &=-\sum_{d\leq2X}\mu(d)\log d
     \sum_{m\geq1}\Lambda(dm+h)W(dm/X)
 \label{eq:exact-source-ledger}\\
 &=\cC_{h,\leq D}(X;W)+\cT_{h,>D}(X;W),
 \label{eq:exact-source-split}
\end{align}
where
\begin{align}
 \cC_{h,\leq D}(X;W)
 &=-\sum_{d\leq D}\mu(d)\log d
     \sum_{m\geq1}\Lambda(dm+h)W(dm/X),
 \label{eq:short-source-ledger}\\
 \cT_{h,>D}(X;W)
 &=-\sum_{D<d\leq2X}\mu(d)\log d
     \sum_{m\geq1}\Lambda(dm+h)W(dm/X).
 \label{eq:large-source-tail}
\end{align}
All these identities are exact.
\end{proposition}

\begin{proof}
Insert \eqref{eq:mobius-log-source-identity} into
\eqref{eq:pair-correlation}, interchange two finite sums, and
write $n=dm$.  Since $W(n/X)$ vanishes unless $X\leq n\leq2X$, no divisor
larger than $2X$ occurs.  Splitting at $D$ proves
\eqref{eq:exact-source-split}.
\end{proof}

Divisors not coprime to $h$ do not form the analytic obstruction.  They are
supported on a finite collection of prime-power targets.

\begin{lemma}[The non-coprime source part]
\label{lem:noncoprime-source-part}
For every $\eps>0$,
\begin{equation}\label{eq:noncoprime-source-bound}
 \sum_{\substack{d\leq2X\\(d,h)>1}}
 |\mu(d)|\log(2d)
 \sum_{m\geq1}\Lambda(dm+h)|W(dm/X)|
 \ll_{h,W,\eps}X^\eps.
\end{equation}
The same bound holds after imposing either $d\leq D$ or $d>D$.
\end{lemma}

\begin{proof}
After writing $n=dm$, the left side is at most
\[
 \sum_{n\asymp X}\Lambda(n+h)|W(n/X)|
 \sum_{\substack{d\mid n\\(d,h)>1}}|\mu(d)|\log(2d).
\]
If the inner sum is nonempty, some prime $p\mid h$ divides $n$.  Hence
$p\mid n+h$.  The condition $\Lambda(n+h)\ne0$ then forces
$n+h=p^k$.  For each of the finitely many primes $p\mid h$, only
$O_h(1)$ powers $p^k$ lie in an interval of the form $[X+h,2X+h]$.
Finally,
\[
 \sum_{d\mid n}|\mu(d)|\log(2d)
 \ll \tau(n)\log(2n)\ll_\eps X^\eps,
\]
after reducing $\eps$ if necessary.  This proves the claim.
\end{proof}

\subsection{The one-sided Bombieri--Vinogradov input}
\label{subsec:one-sided-bv}

We use the maximal form of the Bombieri--Vinogradov theorem.  This is a
classical input, not a new distribution theorem
\citep{Bombieri1965,Vinogradov1965,IwaniecKowalski2004}.

\begin{theorem}[One-sided smooth Type~I estimate]
\label{thm:one-sided-bv-type-I}
Fix $h\ne0$, $A>0$, and $W\in C_c^1((0,\infty))$.  There is a constant
$B=B(A,W)>0$ with the following properties.

Let $x=MN$ be sufficiently large, let $M,N\geq2$, and suppose
\begin{equation}\label{eq:one-sided-bv-range}
 2M\leq \frac{x^{1/2}}{(\log x)^B}.
\end{equation}
For arbitrary complex coefficients $|\alpha_m|\leq1$ supported on
$M<m\leq2M$ and $(m,h)=1$,
\begin{equation}\label{eq:one-sided-bv-row-form}
 \sum_{M<m\leq2M}\alpha_m
 \left\{
   \sum_{n\geq1}\Lambda(mn+h)W(n/N)
   -\frac{mN}{\varphi(m)}\cI_W
 \right\}
 \ll_{A,h,W}\frac{x}{(\log x)^A}.
\end{equation}
In particular, the main term in the $m$th row is
$mN\cI_W/\varphi(m)$, not $N\cI_W$.

There is also a fixed-source-scale form.  If $X$ is sufficiently large,
\begin{equation}\label{eq:fixed-source-bv-range}
 D\leq\frac{X^{1/2}}{(\log X)^B}
\end{equation}
and $|\alpha_d|\leq1$ is supported on $d\leq D$ and $(d,h)=1$, then
\begin{equation}\label{eq:one-sided-bv-fixed-source}
 \sum_{\substack{d\leq D\\(d,h)=1}}\alpha_d
 \left\{
   \sum_{m\geq1}\Lambda(dm+h)W(dm/X)
   -\frac{X}{\varphi(d)}\cI_W
 \right\}
 \ll_{A,h,W}\frac{X}{(\log X)^A}.
\end{equation}
\end{theorem}

\begin{proof}
Put
\[
 E(y;q,a)=
 \sum_{\substack{r\leq y\\r\equiv a\pmod q}}\Lambda(r)
 -\frac{y}{\varphi(q)}
 \qquad ((a,q)=1).
\]
The maximal Bombieri--Vinogradov theorem states that, after increasing
$B$ in terms of $A$, the sum over $q$ up to the right side of
\eqref{eq:one-sided-bv-range} of
$\max_{(a,q)=1}\max_{y\leq Cx}|E(y;q,a)|$ is
$O_A(x(\log x)^{-A})$; here $C$ depends only on the support of $W$.

For fixed $m$, the substitution $r=mn+h$ places the target in the reduced
class $h\pmod m$.  Stieltjes partial summation against
$W((r-h)/(mN))$ shows that the expression in braces in
\eqref{eq:one-sided-bv-row-form} is bounded by a constant depending on
$W$ times
\[
 \max_{y\leq Cx+|h|}|E(y;m,h)|.
\]
The integral of the continuous main term is
\[
 \frac1{\varphi(m)}
 \int_0^\infty W((t-h)/(mN))\,dt
 =\frac{mN}{\varphi(m)}\cI_W.
\]
Summing the error with $|\alpha_m|\leq1$ proves
\eqref{eq:one-sided-bv-row-form}.  The same argument with the weight
$W((r-h)/X)$ gives \eqref{eq:one-sided-bv-fixed-source}; in that case the
continuous main term is $X\cI_W/\varphi(d)$.
\end{proof}

The coefficient $\mu(d)\log d$ in the exact source ledger costs only one
additional logarithm in \cref{thm:one-sided-bv-type-I}.  The main terms are
summed by the following standard singular-series coefficient identity.

\subsection{The singular coefficient and small-divisor evaluation}
\label{subsec:singular-coefficient}

For even $h\ne0$, define
\begin{equation}\label{eq:pair-singular-series}
 \mathfrak S(h)
 =2\prod_{p>2}\left(1-\frac1{(p-1)^2}\right)
   \prod_{\substack{p\mid h\\p>2}}\frac{p-1}{p-2}.
\end{equation}

\begin{lemma}[Singular-series coefficient]
\label{lem:singular-series-coefficient}
For every fixed even $h\ne0$ and every $A>0$,
\begin{equation}\label{eq:singular-coefficient-limit}
 -\sum_{\substack{d\leq z\\(d,h)=1}}
   \frac{\mu(d)\log d}{\varphi(d)}
 =\mathfrak S(h)+O_{A,h}((\log z)^{-A})
 \qquad (z\geq3).
\end{equation}
\end{lemma}

\begin{proof}
For $\Re s>0$, set
\begin{align}
 F_h(s)
 &=\sum_{\substack{d\geq1\\(d,h)=1}}
   \frac{\mu(d)}{\varphi(d)d^s}
 =\prod_{p\nmid h}\left(1-\frac{p^{-s}}{p-1}\right)
 \label{eq:singular-dirichlet-series}\\
 &=\frac{G_h(s)}{\zeta(1+s)},
 \label{eq:singular-zeta-factorization}
\end{align}
where
\begin{equation}\label{eq:singular-G-product}
 G_h(s)=
 \prod_{p\mid h}(1-p^{-1-s})^{-1}
 \prod_{p\nmid h}
 \frac{1-p^{-s}/(p-1)}{1-p^{-1-s}}.
\end{equation}
The second product converges absolutely in a fixed half-plane containing
$s=0$.  At zero its local factors give
\begin{align*}
 G_h(0)
 &=2\prod_{\substack{p>2\\p\nmid h}}
      \frac{p(p-2)}{(p-1)^2}
    \prod_{\substack{p\mid h\\p>2}}\frac{p}{p-1}
 =\mathfrak S(h).
\end{align*}
We now justify the truncation quantitatively rather than differentiate at the
boundary of convergence.  Write
\[
 G_h(s)=\sum_{r\geq1}\frac{g_h(r)}{r^s}.
\]
The coefficients are multiplicative and, for $\nu\geq1$, satisfy
\[
 g_h(p^\nu)=
 \begin{cases}
  p^{-\nu},&p\mid h,\\
  -\{p^\nu(p-1)\}^{-1},&p\nmid h.
 \end{cases}
\]
Hence, for any fixed $0<\eta<1$, the Euler product of the absolute local
sums converges: away from the finitely many primes dividing $h$, its
nonconstant term is $O_\eta(p^{-2+\eta})$.  In particular,
\begin{equation}\label{eq:G-coefficient-moment}
 \sum_{r\geq1}|g_h(r)|r^\eta<\infty.
\end{equation}

Put
\[
 A_0(y)=\sum_{k\leq y}\frac{\mu(k)}k,
 \qquad
 A_1(y)=\sum_{k\leq y}\frac{\mu(k)\log k}{k}.
\]
The classical zero-free region and partial summation imply, for every
$C>0$ and $y\geq1$,
\begin{equation}\label{eq:mobius-over-n-estimates}
 A_0(y)\ll_C(\log(2y))^{-C},
 \qquad
 A_1(y)=-1+O_C((\log(2y))^{-C});
\end{equation}
see \citet{IwaniecKowalski2004,MontgomeryVaughan2007}.
In the common half-plane of absolute convergence,
\eqref{eq:singular-zeta-factorization} gives the coefficient identity
\[
 \frac{\mu(d)\one_{(d,h)=1}}{\varphi(d)}
 =\sum_{rk=d}g_h(r)\frac{\mu(k)}k.
\]
Consequently
\begin{align}
 -\sum_{\substack{d\leq z\\(d,h)=1}}
   \frac{\mu(d)\log d}{\varphi(d)}
 =-\sum_{r\leq z}g_h(r)
 \left\{(\log r)A_0(z/r)+A_1(z/r)\right\}.
 \label{eq:singular-coefficient-convolution}
\end{align}
For $r\leq z^{1/2}$, \eqref{eq:mobius-over-n-estimates}, with $C$
chosen larger than $A+2$, makes the difference between the braces and $-1$
equal to $O_{A,h}((1+\log r)(\log z)^{-A-1})$.  Its sum is admissible by
\eqref{eq:G-coefficient-moment}.  For $r>z^{1/2}$, the elementary bounds
$|A_0(y)|\ll\log(2y)$ and $|A_1(y)|\ll(\log(2y))^2$ combine with
\eqref{eq:G-coefficient-moment} to give
$O_{A,h}((\log z)^{-A})$.  The omitted tail
$\sum_{r>z}g_h(r)$ has the same bound.  Thus
\eqref{eq:singular-coefficient-convolution} equals
$G_h(0)+O_{A,h}((\log z)^{-A})$, and
$G_h(0)=\mathfrak S(h)$ proves the lemma.
\end{proof}

\begin{corollary}[Evaluation of the small source-divisor range]
\label{cor:source-small-divisors}
Fix $A>0$.  There is $B=B(A,h,W)>0$ such that, for all sufficiently
large $X$, with
\begin{equation}\label{eq:small-divisor-cutoff}
 D=\left\lfloor\frac{X^{1/2}}{(\log X)^B}\right\rfloor,
\end{equation}
one has
\begin{equation}\label{eq:short-source-main-term}
 \cC_{h,\leq D}(X;W)
 =\mathfrak S(h)X\cI_W
  +O_{A,h,W}\left(\frac{X}{(\log X)^A}\right).
\end{equation}
\end{corollary}

\begin{proof}
By \cref{lem:noncoprime-source-part}, the terms with $(d,h)>1$ contribute
$O_{h,W,\eps}(X^\eps)$.  On the coprime terms, apply
\eqref{eq:one-sided-bv-fixed-source} with
\[
 \alpha_d=-\frac{\mu(d)\log d}{\log D}.
\]
After multiplying back by $\log D$ and choosing the exponent in the
Bombieri--Vinogradov error larger than $A+2$, this gives
\[
 \cC_{h,\leq D}(X;W)
 =-X\cI_W
  \sum_{\substack{d\leq D\\(d,h)=1}}
  \frac{\mu(d)\log d}{\varphi(d)}
 +O_{A,h,W}(X(\log X)^{-A}).
\]
Now use \cref{lem:singular-series-coefficient}.  Since
$\log D\asymp\log X$, its truncation error is absorbed into the displayed
error.  Taking, for example, $\eps=1/2$ absorbs the non-coprime part as
well.
\end{proof}

The content of \cref{cor:source-small-divisors} is precisely limited
to the short source-divisor range.  Combining it with the exact split gives
an equivalence, not a solution.

\begin{theorem}[Exact large-divisor-tail equivalence]
\label{thm:large-source-tail-equivalence}
For each $A>0$, choose $B=B(A,h,W)$ as in
\cref{cor:source-small-divisors}, and let $D$ be given by
\eqref{eq:small-divisor-cutoff}.  Then, for all sufficiently large $X$,
\begin{equation}\label{eq:pair-minus-main-equals-tail}
 \cC_h(X;W)-\mathfrak S(h)X\cI_W
 =\cT_{h,>D}(X;W)
  +O_{A,h,W}\left(\frac{X}{(\log X)^A}\right).
\end{equation}
Consequently,
\begin{equation}\label{eq:tail-equivalence-asymptotic}
 \cC_h(X;W)\sim\mathfrak S(h)X\cI_W
 \quad\Longleftrightarrow\quad
 \cT_{h,>D}(X;W)=o(X),
\end{equation}
provided $\cI_W\ne0$.
\end{theorem}

\begin{proof}
Substitute \eqref{eq:short-source-main-term} into the exact identity
\eqref{eq:exact-source-split}.  This proves
\eqref{eq:pair-minus-main-equals-tail}, and the equivalence follows.
\end{proof}

\begin{remark}[What the truncation evaluation does not say]
\label{rem:small-divisor-tail-boundary}
The right side of \eqref{eq:pair-minus-main-equals-tail} contains the
divisors $d>D$, where the direct arithmetic-progression argument has no
Bombieri--Vinogradov coverage.  The assertion that this tail is $o(X)$ is
exactly the missing fixed-shift cancellation.  Naming the tail, or noting
that its coefficients have signs, does not estimate it.  In particular,
\cref{thm:large-source-tail-equivalence} is not a reduction of the twin-prime
problem to an independently easier theorem.
\end{remark}

\subsection{A prime-factor/target-prime marginal}
\label{subsec:prime-target-marginal}

The same one-sided theorem gives a standard but distinct prime-target
asymptotic.  Let $W_1,W_2\in C_c^1((0,\infty))$ be supported in $[1,2]$,
put $x=MN$, and define
\begin{align}
 \mathcal M_h(M,N)
 &=\sum_{m,n\geq1}
   \Lambda(m)\Lambda(mn+h)
   W_1(m/M)W_2(n/N),
 \label{eq:prime-target-marginal}\\
 \cP_h(M,N)
 &=\sum_{m,n\geq1}
   \Lambda(m)\Lambda(n)\Lambda(mn+h)
   W_1(m/M)W_2(n/N),
 \label{eq:three-prime-target-form}\\
 \mathcal V_h(M,N)
 &=\sum_{m,n\geq1}
   \Lambda(m)(\Lambda(n)-1)\Lambda(mn+h)
   W_1(m/M)W_2(n/N).
 \label{eq:missing-prime-covariance}
\end{align}

\begin{proposition}[Marginal theorem and covariance ledger]
\label{prop:marginal-covariance-ledger}
Suppose $M,N\to\infty$ and, for a sufficiently large $B$,
\begin{equation}\label{eq:marginal-bv-range}
 2M\leq\frac{x^{1/2}}{(\log x)^B}.
\end{equation}
Then
\begin{equation}\label{eq:prime-target-marginal-asymptotic}
 \mathcal M_h(M,N)
 =MN\left(\int_0^\infty W_1\right)
     \left(\int_0^\infty W_2\right)
 +o(MN).
\end{equation}
Moreover, the identity
\begin{equation}\label{eq:exact-marginal-covariance-ledger}
 \cP_h(M,N)=\mathcal M_h(M,N)+\mathcal V_h(M,N)
\end{equation}
is exact.  Hence the expected three-prime asymptotic, written in a form
that remains meaningful when a weight integral vanishes,
\begin{equation}\label{eq:conjectural-three-prime-main}
 \cP_h(M,N)
 =\mathfrak S(h)MN
 \left(\int_0^\infty W_1\right)
 \left(\int_0^\infty W_2\right)
 +o(MN),
\end{equation}
is equivalent, in this range, to
\begin{equation}\label{eq:required-covariance-main}
 \mathcal V_h(M,N)
 =(\mathfrak S(h)-1)MN
 \left(\int_0^\infty W_1\right)
 \left(\int_0^\infty W_2\right)
 +o(MN).
\end{equation}
In particular, when
$(\mathfrak S(h)-1)(\int W_1)(\int W_2)\ne0$, the missing covariance is not
predicted to be $o(MN)$.
\end{proposition}

\begin{proof}
In \eqref{eq:one-sided-bv-row-form}, take
\[
 C_1=\max\{1,\lVert W_1\rVert_\infty\},
 \qquad
 \alpha_m=\frac{\Lambda(m)W_1(m/M)}{C_1\log(3M)}
\]
and then multiply the result by $C_1\log(3M)$.  Thus
$|\alpha_m|\leq1$ on the support in question.  By choosing the logarithmic
saving one power larger, the resulting error is $o(MN)$.  Rows with
$(m,h)>1$ are negligible: if $\Lambda(m)\ne0$ and a prime $p\mid(m,h)$,
then $m$ is a power of $p$, and $mn+h$ can carry von Mangoldt weight only
when it too is a power of $p$.  For each $p\mid h$, a dyadic interval
contains $O_h(1)$ relevant powers $m=p^a$, and the target interval contains
$O_h(1)$ powers $mn+h=p^b$; their total contribution is
$O_h((\log x)^2)=o(MN)$.

The main term is
\[
 N\left(\int_0^\infty W_2\right)
 \sum_m\Lambda(m)W_1(m/M)\frac{m}{\varphi(m)}.
\]
The prime number theorem and the negligible contribution of higher prime
powers give
\[
 \sum_m\Lambda(m)W_1(m/M)\frac{m}{\varphi(m)}
 =M\int_0^\infty W_1(u)\,du+o(M),
\]
which proves \eqref{eq:prime-target-marginal-asymptotic}.

Equation \eqref{eq:exact-marginal-covariance-ledger} is obtained by writing
$\Lambda(n)=1+(\Lambda(n)-1)$.  For completeness, the local factor of the
three conditions $m,n,mn+h$ is indeed \eqref{eq:pair-singular-series}.  If
$p\nmid h$, there are $(p-1)(p-2)$ pairs $(a,b)\in\mathbb F_p^2$ for which
$a$, $b$, and $ab+h$ are all nonzero.  Relative to three independent prime
conditions, the normalized factor is
\[
 \frac{(p-1)(p-2)/p^2}{(1-1/p)^3}
 =\frac{p(p-2)}{(p-1)^2}.
\]
If $p\mid h$, there are $(p-1)^2$ admissible pairs and the factor is
$p/(p-1)$.  Their product is $\mathfrak S(h)$.  Subtracting the unconditional
marginal asymptotic from \eqref{eq:conjectural-three-prime-main} proves the
equivalence with \eqref{eq:required-covariance-main}.
\end{proof}

\begin{remark}[Two marginals do not determine their intersection]
\label{rem:marginals-do-not-intersect}
The marginal theorem requires $m$ and $mn+h$ to be prime-weighted but leaves
$n$ unrestricted.  The preceding local-kernel results require both factor
variables to be prime-weighted but leave the target only rough.  These two
statements cannot be intersected without estimating
\eqref{eq:missing-prime-covariance}.  Whenever its coefficient is nonzero,
the main term predicted in \eqref{eq:required-covariance-main} quantifies that
missing dependence.
\end{remark}

\subsection{A compact target-side Vaughan ledger}
\label{subsec:compact-vaughan-ledger}

For later use, we finish with an exact target-side decomposition.  It is a
standard form of Vaughan's identity
\citep{Vaughan1977,IwaniecKowalski2004}; generalized identities produce more
multilinear packets but do not estimate them by themselves
\citep{HeathBrownVaughan1982}.

\begin{proposition}[Exact Vaughan ledger]
\label{prop:exact-vaughan-ledger}
Let $F$ be any finitely supported function on the positive integers and let
$U,V\geq1$.  Put
\begin{equation}\label{eq:vaughan-beta-coefficient}
 \beta_V(k)=\sum_{\substack{d\mid k\\d\leq V}}\mu(d).
\end{equation}
Then
\begin{align}
 \sum_t\Lambda(t)F(t)
={}&\sum_{t\leq U}\Lambda(t)F(t)
 -\sum_{\substack{\ell\leq U\\d\leq V}}
   \Lambda(\ell)\mu(d)\sum_{r\geq1}F(\ell dr)
 \notag\\
 &+\sum_{d\leq V}\mu(d)\sum_{r\geq1}(\log r)F(dr)
 -\sum_{\substack{\ell>U\\k>1}}
   \Lambda(\ell)\beta_V(k)F(\ell k).
 \label{eq:exact-vaughan-four-packets}
\end{align}
Moreover, $\beta_V(k)=0$ for $1<k\leq V$.
\end{proposition}

\begin{proof}
Write $\mu_{\leq V}(d)=\mu(d)\one_{d\leq V}$, so that
$\beta_V=\mu_{\leq V}*\one$.  The second term in
\eqref{eq:exact-vaughan-four-packets} is
$-(\Lambda_{\leq U}*\beta_V)(t)$, while the third is
$(\mu_{\leq V}*\log)(t)$.  The last term equals
\[
 - (\Lambda_{>U}*\beta_V)(t)+\Lambda_{>U}(t),
\]
because the omitted value $k=1$ has $\beta_V(1)=1$.  Adding the first term
therefore gives $\Lambda(t)$ plus
\[
 (\mu_{\leq V}*\log)(t)-(\Lambda*\mu_{\leq V}*\one)(t),
\]
which is zero because $\Lambda*\one=\log$.  Finally, if
$1<k\leq V$, then every divisor of $k$ is at most $V$, and
$\beta_V(k)=\sum_{d\mid k}\mu(d)=0$.
\end{proof}

To see the remaining geometry, take
\begin{equation}\label{eq:vaughan-factor-incidence}
 F_h^{a,b}(t)=\sum_{m,n}a_mb_n\one_{mn+h=t}.
\end{equation}
Here the coefficients are finite and satisfy
$a_mb_n\ne0\Rightarrow mn+h\geq1$.
The middle two packets in \eqref{eq:exact-vaughan-four-packets} become the
Type~I congruence sums
\begin{align}
 &-\sum_{\substack{\ell\leq U\\d\leq V}}
  \Lambda(\ell)\mu(d)
  \sum_{mn\equiv-h\, (\mathrm{mod}\,\ell d)}a_mb_n,
 \label{eq:vaughan-type-I-first}\\
 &\sum_{d\leq V}\mu(d)
  \sum_{mn\equiv-h\, (\mathrm{mod}\,d)}
  a_mb_n\log\left(\frac{mn+h}{d}\right),
 \label{eq:vaughan-type-I-second}
\end{align}
whereas the last packet is
\begin{equation}\label{eq:vaughan-hard-packet}
 -\sum_{\substack{\ell>U\\k>V}}
  \Lambda(\ell)\beta_V(k)
  \sum_{mn+h=\ell k}a_mb_n.
\end{equation}
Thus the unresolved fixed-shift object is supported on the shifted
multiplicative surface
\begin{equation}\label{eq:vaughan-shifted-surface}
 \ell k-mn=h,
 \qquad \ell>U,\quad k>V.
\end{equation}
The finite roughness kernel tests congruence avoidance by prescribed primes;
it does not estimate \eqref{eq:vaughan-hard-packet}.  Consequently the exact
Vaughan ledger identifies a Type~II interface but does not close it.
